\chapter{Estado del arte}
\label{cap:estado-arte}

% Objetivo: ~2 páginas. Qué está pasando con los agentes LLM y la IA en mercados,
% literatura sobre supervisión en tiempo de ejecución, y el hueco que llena este TFG.
% Citas OBLIGATORIAS (revisor-bibliografico buscará los papers académicos del estado del arte).

% [PENDIENTE DE REDACCIÓN — pipeline /redaccion-capitulo]
% Estructura prevista:
%   - Agentes LLM aplicados a decisiones financieras (estado actual, limitaciones).
%   - Detección de anomalías / supervisión runtime de sistemas automáticos.
%   - Detección de regímenes y gestión de riesgo clásica como capa de control.
%   - El hueco: supervisión estadística interpretable sobre un agente LLM. Aquí entra STRATA.
